\documentclass[11pt,a4paper]{article}
\usepackage[margin=1in]{geometry}
\usepackage{hyperref}
\usepackage{listings}
\usepackage{xcolor}
\usepackage{longtable}
\usepackage{booktabs}
\usepackage{fancyhdr}

\hypersetup{colorlinks=true,linkcolor=blue,urlcolor=blue}

\lstset{
  basicstyle=\ttfamily\small,
  breaklines=true,
  frame=single,
  backgroundcolor=\color{gray!10},
  keywordstyle=\color{blue},
  commentstyle=\color{green!50!black},
  stringstyle=\color{red!70!black},
}

\pagestyle{fancy}
\fancyhead[L]{sheLLaMa Admin Manual}
\fancyhead[R]{\thepage}
\fancyfoot{}

\title{\textbf{sheLLaMa Admin Manual}\\[0.5em]\large Deployment, Configuration \& Operations}
\author{Version 1.0 --- May 2026}
\date{}

\begin{document}
\maketitle
\tableofcontents
\newpage

%======================================================================
\section{Architecture Overview}

\begin{verbatim}
Clients --> Frontend (:5000) --> Backend Farm
                                  |-- Backend 1 (:5000)
                                  |-- Backend 2 (:5000)
                                  |-- Backend N (:5000)
External tools --> /v1/chat/completions (OpenAI-compatible)
\end{verbatim}

\begin{itemize}
  \item \textbf{Backend} (\texttt{backend/app.py}) --- Ollama worker, task queue, cloud fallback, image generation subprocess
  \item \textbf{Frontend} (\texttt{frontend/app-distributed.py}) --- Load balancer, routing, caching, auth, rate limiting, health checks, webhooks, admin console
  \item \textbf{Clients} --- Bash CLI, PowerShell CLI, integrations, GUIs, REST API consumers
\end{itemize}

%======================================================================
\section{System Requirements}

\subsection{Backend Server}
\begin{longtable}{p{3cm}p{2cm}p{2cm}p{2cm}p{5cm}}
\toprule
\textbf{Tier} & \textbf{CPU} & \textbf{RAM} & \textbf{Storage} & \textbf{Models} \\
\midrule
Minimum & 8 cores & 16GB & 50GB & 7B (30--60s/response) \\
Recommended & 16 cores & 32GB & 100GB & 13B--14B (1--3min/response) \\
Large & 32+ cores & 64GB+ & 200GB & 32B+ (not recommended for CPU) \\
\bottomrule
\end{longtable}

\subsection{Frontend Server}
4 cores, 8GB RAM, 20GB storage.

%======================================================================
\section{Deployment}

\subsection{Quick Start (Single Server)}
\begin{lstlisting}[language=bash]
# Install Ollama
curl -fsSL https://ollama.com/install.sh | sh
ollama pull qwen2.5-coder:7b

# Clone and run
git clone https://github.com/pfworks/sheLLaMa.git
cd shellama
python3 -m venv venv && source venv/bin/activate
pip install flask ollama pyyaml requests psutil
python backend/app.py
\end{lstlisting}

\subsection{Distributed Deployment (Ansible)}
\begin{lstlisting}[language=bash]
cp deploy/inventory.ini.example inventory.ini
cp deploy/inventory-frontend.ini.example inventory-frontend.ini
cp deploy/backends.json.example backends.json
# Edit each file with your server details

ansible-playbook -i inventory.ini deploy/deploy.yml
ansible-playbook -i inventory-frontend.ini deploy/deploy-frontend.yml
\end{lstlisting}

\subsection{Manual Deployment}
\begin{lstlisting}[language=bash]
# Frontend
scp frontend/app-distributed.py root@frontend:/usr/local/bin/
scp shared/*.py root@frontend:/usr/local/shared/
scp frontend/web/*.html root@frontend:/export/html/
ssh root@frontend "systemctl restart ansible-ollama-frontend"

# Backend
scp backend/app.py root@backend:/usr/local/bin/
scp shared/*.py root@backend:/usr/local/shared/
ssh root@backend "systemctl restart ansible-ollama"
\end{lstlisting}

\subsection{backends.json}
\begin{lstlisting}[language=bash]
{
  "backends": [
    {"url": "http://192.168.1.230:5000", "weight": 10,
     "max_model": "qwen2.5-coder:7b"},
    {"url": "http://192.168.1.233:5000", "weight": 10,
     "max_model": "qwen2.5-coder:14b"}
  ]
}
\end{lstlisting}
Higher weight = higher priority. Score = \texttt{queue\_size - (weight * 0.1)}, lowest wins.

%======================================================================
\section{Admin Console}

Access at \texttt{http://frontend:5000} (redirects to \texttt{/status}).

\begin{longtable}{p{3cm}p{3cm}p{8cm}}
\toprule
\textbf{Page} & \textbf{URL} & \textbf{Description} \\
\midrule
Status & \texttt{/status} & Total requests, tokens, active backends, queue, fallback toggle \\
Backends & \texttt{/backends} & Per-backend: online/offline, CPU/RAM, weight, models, active task \\
Stats & \texttt{/stats} & Charts: queue size and token usage over time \\
Costs & \texttt{/costs} & Cloud cost tracking by time range, fallback spend \\
Cloud Usage & \texttt{/cloud-usage} & Per-IP and per-API-key cost breakdown \\
Settings & \texttt{/settings} & Model aliases, toggles, webhooks, command permissions \\
\bottomrule
\end{longtable}

%======================================================================
\section{Command Permissions}

The agentic shell uses three permission tiers for commands:

\begin{longtable}{p{3cm}p{3cm}p{7.5cm}}
\toprule
\textbf{Tier} & \textbf{Behavior} & \textbf{Examples} \\
\midrule
Allow & Auto-run & ls, cat, grep, git status, curl, ps, echo \\
Prompt & Asks user & pip install, python3, mkdir, git commit \\
Block & Always refused & rm -rf /, mkfs, git push --force, fork bombs \\
\bottomrule
\end{longtable}

\subsection{Managing Permissions}

\textbf{Via Settings page:} Navigate to \texttt{/settings} $\rightarrow$ Command Permissions.
\begin{itemize}
  \item \textbf{Auto-Allowed Patterns}: editable regex list. Add/remove freely.
  \item \textbf{Blocked --- Built-in}: immutable safety rules. Cannot be removed.
  \item \textbf{Blocked --- Custom}: admin-added block rules. Can be removed.
\end{itemize}

\textbf{Via API:}
\begin{lstlisting}[language=bash]
# Get current permissions
curl http://frontend:5000/api/permissions

# Add an allow pattern
curl -X POST http://frontend:5000/api/permissions \
  -H "Content-Type: application/json" \
  -d '{"add_allowed": "^npm\\s"}'

# Add a custom block pattern
curl -X POST http://frontend:5000/api/permissions \
  -H "Content-Type: application/json" \
  -d '{"add_blocked": "dangerous_script\\.sh"}'

# Remove a custom block (built-in blocks cannot be removed)
curl -X POST http://frontend:5000/api/permissions \
  -H "Content-Type: application/json" \
  -d '{"remove_blocked": "dangerous_script\\.sh"}'
\end{lstlisting}

The CLI fetches permissions from the server on startup (2s timeout, falls back to hardcoded defaults if unreachable).

%======================================================================
\section{Authentication}

Optional---disabled when \texttt{/etc/shellama/auth.json} doesn't exist.

\subsection{API Keys}
\begin{lstlisting}[language=bash]
# Client usage
export SHELLAMA_API_KEY=sk-your-key
# Or header: X-API-Key: sk-your-key
\end{lstlisting}

\subsection{SSO (OIDC)}
Supports Keycloak, Azure AD, Authentik. Configure in \texttt{auth.json}:
\begin{lstlisting}[language=bash]
cp deploy/auth.json.example /etc/shellama/auth.json
# Edit with your SSO provider details
\end{lstlisting}

\subsection{Roles}
\begin{longtable}{p{2cm}p{3cm}p{3cm}p{3cm}p{2.5cm}}
\toprule
\textbf{Role} & \textbf{API Access} & \textbf{Web UI} & \textbf{Models} & \textbf{Cloud} \\
\midrule
admin & All endpoints & Full (modify) & All & Yes \\
user & Chat, gen, explain & View only & Configured & Configurable \\
viewer & Read-only & View only & None & No \\
\bottomrule
\end{longtable}

%======================================================================
\section{Cloud Fallback}

When local models produce poor output, requests can automatically retry via cloud.

\subsection{OpenRouter (Cloud)}
\begin{lstlisting}[language=bash]
# On each backend (systemd override):
Environment="OPENROUTER_API_KEY=sk-or-v1-your-key"
Environment="OPENROUTER_MODEL=anthropic/claude-3.5-sonnet"
Environment="USE_CLOUD_FALLBACK=true"
\end{lstlisting}

\subsection{LiteLLM (Self-Hosted)}
\begin{lstlisting}[language=bash]
Environment="OPENROUTER_API_KEY=sk-anything"
Environment="OPENROUTER_MODEL=fallback"
Environment="OPENROUTER_URL=http://litellm-host:4000/v1/chat/completions"
Environment="USE_CLOUD_FALLBACK=true"
\end{lstlisting}

\subsection{Auto-Fallback Mode}
Toggle via Status page or API:
\begin{lstlisting}[language=bash]
curl -X POST http://frontend:5000/auto-fallback \
  -H "Content-Type: application/json" \
  -d '{"enabled": true}'
\end{lstlisting}
When enabled, poor local responses auto-retry via cloud without client confirmation.

%======================================================================
\section{Load Balancing}

\begin{itemize}
  \item Weighted routing: score = \texttt{queue\_size - (weight * 0.1)}, lowest wins
  \item Model size filtering: requested model must be $\leq$ backend's \texttt{max\_model}
  \item Multi-file analysis runs in parallel across backends
  \item Health checks every 30s; 3 consecutive failures $\rightarrow$ unhealthy (skipped)
  \item Auto-recovery when backend responds again
  \item Retry up to 2x on different backends on failure
\end{itemize}

%======================================================================
\section{Rate Limiting \& Budgets}

Per-key configuration in \texttt{auth.json}:
\begin{lstlisting}[language=bash]
"api_keys": {
  "sk-team-key": {
    "name": "team",
    "role": "user",
    "rate_limit": {"rpm": 30, "tpd": 500000},
    "budget": {"max_daily": 5.00}
  }
}
\end{lstlisting}

\begin{itemize}
  \item \texttt{rpm}: requests per minute
  \item \texttt{tpd}: tokens per day
  \item \texttt{budget.max\_daily}: max cloud fallback spend per day (USD)
  \item Budget warning webhook fires at 80\%
  \item Returns HTTP 429 when exceeded
\end{itemize}

%======================================================================
\section{Prompt Caching}

\begin{itemize}
  \item Key: SHA256(endpoint + model + content)
  \item TTL: 5 minutes (configurable via \texttt{SHELLAMA\_CACHE\_TTL}, 0 to disable)
  \item Max entries: 500
  \item Skips: conversations, force\_cloud, errors
  \item Stats tracked: cached\_requests, tokens\_saved
\end{itemize}

%======================================================================
\section{OpenAI-Compatible API}

Use sheLLaMa as a drop-in replacement for OpenAI:
\begin{lstlisting}[language=bash]
export OPENAI_API_BASE=http://frontend:5000/v1
export OPENAI_API_KEY=sk-your-key

curl -X POST http://frontend:5000/v1/chat/completions \
  -H "Content-Type: application/json" \
  -H "Authorization: Bearer sk-your-key" \
  -d '{"model": "fast", "messages": [{"role": "user",
       "content": "Hello"}]}'
\end{lstlisting}

Works with: Cursor, Continue, Open WebUI, LangChain, any OpenAI client. Model aliases supported (\texttt{fast} $\rightarrow$ \texttt{llama3.2:1b}).

%======================================================================
\section{TLS \& mTLS}

\subsection{Certificate Management}
\begin{lstlisting}[language=bash]
bin/generate-certs.sh init                    # Generate CA
bin/generate-certs.sh server backend-1 \
  192.168.1.230,localhost                     # Server cert
bin/generate-certs.sh client frontend-mtls   # Client cert
bin/generate-certs.sh list                   # List all certs
bin/generate-certs.sh revoke backend-1       # Revoke
\end{lstlisting}

\subsection{Environment Variables}
\begin{longtable}{p{5cm}p{8.5cm}}
\toprule
\textbf{Variable} & \textbf{Description} \\
\midrule
\texttt{SHELLAMA\_TLS\_CERT} & Server TLS certificate (enables HTTPS) \\
\texttt{SHELLAMA\_TLS\_KEY} & Server TLS private key \\
\texttt{SHELLAMA\_TLS\_CA} & CA cert for client verification (mTLS) \\
\texttt{SHELLAMA\_BACKEND\_CERT} & Client cert for frontend$\rightarrow$backend mTLS \\
\texttt{SHELLAMA\_BACKEND\_KEY} & Client key for frontend$\rightarrow$backend mTLS \\
\texttt{SHELLAMA\_BACKEND\_CA} & CA to verify backend server certs \\
\bottomrule
\end{longtable}

%======================================================================
\section{Webhooks}

Events: \texttt{backend\_down}, \texttt{backend\_recovered}, \texttt{budget\_warning}.

Configure via Settings page or API:
\begin{lstlisting}[language=bash]
curl -X POST http://frontend:5000/api/webhooks \
  -H "Content-Type: application/json" \
  -d '{"add": "https://hooks.example.com/shellama"}'
\end{lstlisting}

Same event suppressed for 5 minutes (dedup).

%======================================================================
\section{Cloud Cost Estimation}

sheLLaMa tracks every token and projects costs across 28+ cloud models.

\subsection{Costs Page (/costs)}
\begin{itemize}
  \item Time range filtering: day/week/month/year/custom
  \item Hypothetical costs: what local usage would cost on cloud
  \item Actual fallback spend: real costs from cloud fallback
  \item Cached savings: tokens served from cache
  \item Cheapest/costliest provider highlighted
\end{itemize}

\subsection{Cloud Usage Page (/cloud-usage)}
\begin{itemize}
  \item Per-IP cost breakdown with expandable provider detail
  \item Per-API-key cost breakdown
  \item Per-cloud-model usage (fallback only)
  \item Cost range (cheapest to most expensive) per client
\end{itemize}

\subsection{Enterprise Comparison}
\begin{lstlisting}[language=bash]
curl "http://frontend:5000/enterprise-costs?users=10&days=30"
\end{lstlisting}
Compares seat-based subscriptions (Copilot, ChatGPT Team, Claude Team, etc.) vs equivalent per-token API spend.

%======================================================================
\section{Monitoring \& Operations}

\subsection{Service Management (Linux)}
\begin{lstlisting}[language=bash]
sudo systemctl status ansible-ollama          # Backend
sudo systemctl status ansible-ollama-frontend # Frontend
sudo journalctl -u ansible-ollama -f          # Logs
sudo systemctl restart ansible-ollama
\end{lstlisting}

\subsection{Health Checks}
\begin{lstlisting}[language=bash]
curl http://frontend:5000/queue-status   # Aggregate status
curl http://backend:5000/queue-status    # Single backend
\end{lstlisting}

\subsection{Reset Counters}
\begin{lstlisting}[language=bash]
curl -X POST http://frontend:5000/reset-stats       # Requests/tokens
curl -X POST http://frontend:5000/reset-cloud-costs # Cost data
curl -X POST http://frontend:5000/reset-all         # Everything
\end{lstlisting}

\subsection{Task Timeouts}
\begin{itemize}
  \item Default: 1800s (30 min) per task (\texttt{SHELLAMA\_TASK\_TIMEOUT})
  \item Stale task reaper checks every 10s
  \item Tasks without heartbeat for 30s are killed
  \item Frontend$\rightarrow$backend timeout: 3600s (1 hour)
\end{itemize}

%======================================================================
\section{Environment Variables (Complete)}

\begin{longtable}{p{5.5cm}p{3.5cm}p{5cm}}
\toprule
\textbf{Variable} & \textbf{Default} & \textbf{Description} \\
\midrule
\texttt{SHELLAMA\_API} & 192.168.1.229:5000 & API endpoint (clients) \\
\texttt{SHELLAMA\_MODEL} & auto & Default model \\
\texttt{SHELLAMA\_API\_KEY} & (empty) & API key \\
\texttt{SHELLAMA\_TASK\_TIMEOUT} & 1800 & Max task runtime (backend) \\
\texttt{SHELLAMA\_COMPACT\_THRESHOLD} & 24000 & Auto-compact threshold \\
\texttt{SHELLAMA\_CACHE\_TTL} & 300 & Prompt cache TTL (0=off) \\
\texttt{SHELLAMA\_AUTH\_FILE} & /etc/shellama/auth.json & Auth config \\
\texttt{SHELLAMA\_WEBHOOK\_URL} & (empty) & Webhook URL \\
\texttt{SHELLAMA\_CERT\_DIR} & /etc/shellama/pki & PKI directory \\
\texttt{USE\_CLOUD\_FALLBACK} & false & Enable cloud fallback \\
\texttt{OPENROUTER\_API\_KEY} & (empty) & Cloud fallback key \\
\texttt{OPENROUTER\_MODEL} & claude-3.5-sonnet & Cloud fallback model \\
\texttt{OPENROUTER\_URL} & openrouter.ai/... & Cloud endpoint \\
\texttt{AI\_IMAGE\_MODEL} & sdxl-turbo & Image model \\
\texttt{SHELLAMA\_DOWNLOAD\_DIR} & (current dir) & Image save dir \\
\bottomrule
\end{longtable}

%======================================================================
\section{Troubleshooting}

\textbf{No backends available} --- Check \texttt{backends.json}; \texttt{max\_model} must match or exceed the requested model.

\textbf{Timeouts} --- Frontend$\rightarrow$backend timeout is 3600s. If tasks exceed \texttt{SHELLAMA\_TASK\_TIMEOUT}, the reaper kills them.

\textbf{Slow responses} --- Use smaller models. Add more backends. Check \texttt{htop}.

\textbf{Image generation slow} --- Use \texttt{sd-turbo} instead of \texttt{sdxl-turbo}. First request loads model (\~{}90s), subsequent are faster.

\textbf{Permission endpoint returns 403} --- Cannot remove built-in safety rules. Only custom block rules can be removed.

\textbf{Cloud fallback not triggering} --- Verify \texttt{USE\_CLOUD\_FALLBACK=true} on backends. Check \texttt{OPENROUTER\_API\_KEY} is set.

\end{document}
